% BOOK-FILE-ID: f-0ade19ade1d2
\subsection{ニュートン法}\label{節：ニュートン法}

曲線の接線を使って,方程式の解の近似値を順に求める方法を考えよう.\footnote{微分が未習なら,接線と$x$軸の交点から次の近似値を作る図形的な発想と,得られた\bookterm{recurrence}{漸化式}の変形を先に読んでよい.}
曲線と$x$軸の交点を直接求める代わりに,今いる点の近くで曲線を接線に置き換える.
直線の交点なら一次方程式で求められる.その位置を次の近似値にし,同じ操作を繰り返す.
$x=a$の近くでは,
\[
  f(a+h) \approx f(a)+f'(a)h
\]
と近似できる.方程式
\[
  f(x)=x^2-2=0
\]
の正の解を求めよう.正の近似値$a_n$での接線は,傾き$2a_n$より,
\[
  y=f(a_n)+2a_n(x-a_n)
\]
である.\,$y=0$ と置いて $x$ 軸との交点を求めると,
\begin{align*}
  a_{n+1}
  &=a_n-\frac{f(a_n)}{f'(a_n)}\\
  &=a_n-\frac{a_n^2-2}{2a_n}\\
  &=\frac{a_n^2+2}{2a_n}
\end{align*}
を得る.これが\bookterm{newton}{ニュートン法}の\bookterm{recurrence}{漸化式}である.
$a_1>0$から各項も正となり,分母は0にならない.
\BookFigNewton

例えば $a_1=2$ とすると,
\[
  a_2=\frac{3}{2},
  \quad
  a_3=\frac{17}{12},
  \quad
  a_4=\frac{577}{408}
\]
となり,これらは $\sqrt{2}$ に急速に近づく.この速さは,誤差を
$e_n=a_n-\sqrt{2}$ と置くと見える.
\begin{align*}
  e_{n+1}
  &=\frac{a_n^2+2}{2a_n}-\sqrt{2}\\
  &=\frac{(a_n-\sqrt{2})^2}{2a_n}\\
  &=\frac{e_n^2}{2a_n}
\end{align*}
誤差が小さくなると,分母は$2\sqrt2$に近づき,誤差は二乗に応じて減る.
\bookterm{general}{一般項}を求めるには,残る分母$a_n$を消したい.
$a_n-\sqrt2$と$a_n+\sqrt2$を並べると,
\BookMathLayout{
  \[
    \begin{aligned}
      a_{n+1}-\sqrt2&=\frac{(a_n-\sqrt2)^2}{2a_n},\\
      a_{n+1}+\sqrt2&=\frac{(a_n+\sqrt2)^2}{2a_n}
    \end{aligned}
  \]
}{
  \[
    \begin{aligned}
      a_{n+1}-\sqrt2&=\frac{(a_n-\sqrt2)^2}{2a_n},\\
      a_{n+1}+\sqrt2&=\frac{(a_n+\sqrt2)^2}{2a_n}
    \end{aligned}
  \]
}
となる.2つの分母が同じなので,比を取れば$2a_n$が消えてくれる.
\[
  \frac{a_{n+1}-\sqrt{2}}{a_{n+1}+\sqrt{2}}
  =\left(\frac{a_n-\sqrt{2}}{a_n+\sqrt{2}}\right)^{2}
\]
となる.そこで
\[
  b_n=\frac{a_n-\sqrt{2}}{a_n+\sqrt{2}}
\]
と置けば,\,$b_{n+1}=b_n^2$となる.これなら順に二乗して,\,$b_n=b_1^{2^{n-1}}$と求められる.正の$b_1$の場合は,\sectionref{節：対数型}と同じく対数を取ってもよい.
そのまま2乗されるのは誤差$e_n$ではなく,比$b_n$の方なのだ.
この置き換えに見覚えはないだろうか.\,$\sqrt2$と$-\sqrt2$は漸化式の2つの不動点であり,\,$b_n$は\sectionref*{節：一次分数型}で使った「2つの不動点からのずれの比」そのものである.

\begin{bookpoint}
  \textbf{POINT}\par
  $\sqrt2$を求めるこのニュートン法では,2つの不動点からのずれの比を取ると,分母が消えて$b_{n+1}=b_n^2$の形になる.
\end{bookpoint}

$a_1=2$のとき,\,$b_1=\frac{2-\sqrt2}{2+\sqrt2}=3-2\sqrt2=\frac1{3+2\sqrt2}$であるから,\,$u=3+2\sqrt2$とおくと$b_n=u^{-2^{n-1}}$である.これを$b_n=\frac{a_n-\sqrt2}{a_n+\sqrt2}$について$a_n$で解けば,
\[
  a_n=\sqrt{2}\,\frac{u^{2^{n-1}}+1}{u^{2^{n-1}}-1},\qquad u=3+2\sqrt2
\]
という,\,$n$から直接計算できる式が得られる.
$u>1$なので,\,$u^{-2^{n-1}}$は$0$に近づき,\,$a_n$は$\sqrt2$に近づくことも分かる.

この式は$f(x)=x^2-2$に固有の変形から得たものである.
一般の\bookterm{newton}{ニュートン法}でも得られるとは限らないが,収束先や誤差は調べられる.
\sectionref*{節：和と差から微積分へ}では,このように漸化式を作る考え方を収束,\bookterm{fixed}{不動点},誤差評価,微分方程式へ広げる.
